\section{Adversarial Reframing}
\label{sec:adversarial}

\subsection{Why Parameters Might Be Partisan}

Before measuring the actual partisan effect of parameter tuning, it is
worth asking: \emph{why might parameters be partisan at all?}
If the algorithm cannot see partisan data, how could parameter choices
produce partisan outcomes?

The answer lies in the correlation between parameter-sensitive outcomes
and political geography.
Parameters that affect compactness indirectly affect partisanship because
urban areas (which vote heavily Democratic) are geographically compact.
A parameter that increases compactness also increases the tendency to
pack Democratic voters into fewer districts.
A parameter that decreases compactness does the opposite.

This correlation suggests a mechanism: by setting parameters that produce
maximally compact districts, a Republican-preferring actor might increase
the packing of Democrats and produce a more Republican-favorable map.
Conversely, a Democrat-preferring actor might set parameters that reduce
compactness and spread Democratic voters across more districts.

\subsection{The Geographic Continuity Argument}

U.2~\citep{b17paper} identifies the structural reason why this mechanism
is weaker than intuition suggests: the bisection tree structure is determined
by the prime factorisation of the district count $k$ and the graph topology,
not by parameter values.
Parameter changes alter \emph{leaf-level cut quality} within a fixed tree
structure, not the tree structure itself.

In practice, the difference between a compact leaf-level cut and a less
compact one involves moving $1$--$3$ census tracts across a district
boundary.
Competitive-district geography---the geography of swing districts that
determine which party controls the delegation---is not sensitive to
such small boundary shifts.
The tracts that move across boundaries are typically in non-competitive
regions of the state, where a shift does not change the partisan winner.

\subsection{Formal Decomposition}

Let $D_i(\theta)$ denote the Democratic vote share in district $i$ under
parameter vector $\theta = (\ufactor, \acounty, T, \aswing, \wvra)$.
The expected number of Democratic seats nationally is:
\[
  \Dnat(\theta) = \sum_{s=1}^{50} \sum_{i \in s} \mathbf{1}[D_i(\theta) > 0.5].
\]

A parameter change $\theta \to \theta'$ shifts $\Dnat$ by:
\[
  \Delta \Dnat = \sum_{s=1}^{50} \sum_{i \in s}
  \left( \mathbf{1}[D_i(\theta') > 0.5] - \mathbf{1}[D_i(\theta) > 0.5] \right).
\]

A district $i$ changes partisan outcome only if $D_i(\theta)$ crosses 0.5
between $\theta$ and $\theta'$.
This requires that district $i$ is competitive (close to the 50\% threshold)
and that the parameter change moves boundary tracts with partisan
composition different from the district's current majority.
Both conditions must hold simultaneously.

The empirical rarity of this conjunction---competitive districts whose
boundaries shift across parameters AND whose shifting tracts contain
enough partisan votes to flip the district---explains the 0.3-seat bound.
